\documentclass[onecolumn,english,12pt]{IEEEtran}
\usepackage[T1]{fontenc}
\usepackage{babel}
\usepackage{amsmath}
\usepackage{amsthm}
\usepackage{amsfonts}
\usepackage{graphicx}
\usepackage{booktabs}
\usepackage[caption=false,font=footnotesize]{subfig}
\usepackage[section]{placeins}

\makeatletter
\theoremstyle{plain}
\newtheorem{thm}{\protect\theoremname}
\theoremstyle{plain}
\newtheorem{lem}[thm]{\protect\lemmaname}
\theoremstyle{plain}
\newtheorem{cor}[thm]{\protect\corollaryname}
\newtheorem{prop}{\protect\propositionname}
\newtheorem{claim}{\protect\claimname}
\newtheorem{conjecture}{\protect\conjecturename}
\theoremstyle{definition}
\newtheorem{example}[thm]{\protect\examplename}
\makeatother

\providecommand{\corollaryname}{Corollary}
\providecommand{\lemmaname}{Lemma}
\providecommand{\theoremname}{Theorem}
\providecommand{\examplename}{Example}
\providecommand{\propositionname}{Proposition}
\providecommand{\claimname}{Claim}
\providecommand{\conjecturename}{Conjecture}


\begin{document}

%% ============================================================
%%  EE 496 CAPSTONE COVER SHEET
%% ============================================================
\begin{titlepage}
  \normalfont

  \textit{Department of Electrical Engineering, University of Hawaii, \today}

  \vspace{0.5em}
  \textit{EE 496 Final Report}

  \vspace{4cm}

  {\Huge Transformer Based Lossless Source Coding}

  \vspace{3cm}

  {\large EE 496 Final Report}

  {\large Spring 2026}

  \vspace{1.5em}

  {\large Department of Electrical Engineering}

  {\large University of Hawaii at M\={a}noa}

  \vspace{2cm}

  {\Large Dan Vu}

  \vspace{3cm}

  \today

  \vfill

  {\large Faculty Advisor: Professor Anders H\o{}st-Madsen}

  {\large Support: Yixin Zhang}

\end{titlepage}

%% ============================================================
%%  PAPER BEGINS HERE (IEEEtran format, same as Madsen's paper)
%% ============================================================

\author{Dan Vu%
  \thanks{Department of Electrical Engineering,
          University of Hawaii at M\={a}noa, Honolulu, HI 96822.
          Advised by Prof.\ Anders H\o{}st-Madsen.}}

\begin{abstract}
H{\o}st-Madsen et al.~\cite{HostMadsen24} discussed that a learned source coder's
redundancy decays as $1/m$ in the number of training sequences, demonstrating that
offline training can beat universal online coders once $m$ is large enough.
This paper empirically tests that prediction using a small transformer
(minGPT-nano, $\sim$0.11M parameters) as a sequential probability estimator
coupled with arithmetic coding, benchmarked against the Laplace estimator and
Context Tree Weighting (CTW) across IID binary, Markov, and FSM sources.
On IID sources, a general transformer trained across the full source distribution
recovers the Laplace estimator from data alone. A specialized transformer
crosses below the Laplace redundancy at around $m = 100$ training sequences.
On structured sources the advantage of offline training widens. The specialized
transformer beats both Laplace and CTW at all sequence lengths on the FSM
source, with CTW closing the gap steadily.
The results emphasize that learned and universal compression operate on the same
mathematical engine of sequential prediction, and that offline knowledge of the
source distribution is the mechanism by which a learned coder earns its advantage.
\end{abstract}

%% ------------------------------------------------------------
\section{Introduction}
\label{sec:intro}

\subsection{Why Compression Matters}
\label{sec:why}

Compression comes in two flavors. Lossy compression discards information
the receiver cannot perceive or does not need. JPEG drops high-frequency
detail and MP3 removes inaudible components. Lossless compression discards
nothing. Every binary bit of the original file can be recovered exactly. Lossless compression
is required wherever one wrong bit corrupts the output. Some examples include executable code,
legal records, medical images, cryptographic keys. There is no tolerance
for approximation in those domains.

It turns out that the problem of lossless compression reduces to modeling the source distribution. A
compressor assigns short codes to common patterns and long codes to rare
ones. The better the probability model matches the true distribution,
the shorter the average code. At scale this compounds: streaming
platforms compress billions of files and digital archives grow by
terabytes per day. A marginal improvement in the probability model
shortens every symbol in every file. In this paper, we demonstrate the use of transformers as compressors, showing that data prediction is the same mathematical engine as data compression.
\subsection{Entropy and Optimal Codes}
\label{sec:entropy}

In this paper, we test our compression models with IID data. We test with binary because it is a simple alphabet that can produce very specialized source distributions. Additionally, all modern digital data can be represented in binary.

A source emits binary symbols
$x_1, x_2, \ldots$ drawn from $\{0, 1\}$ from a distribution $p$. A code assigns
a binary string to each symbol. We find an optimal coding scheme that minimizes the expected
codelength $\mathbb{E}[L(X)]$. The optimal codelength assigned to symbol $i$ which has a probability $p_i$ is: 

\begin{equation}
  l_i^* = -\log_2 p_i.
  \label{eq:optlen}
\end{equation}
The expected optimal codelength is therefore
\begin{equation}
  H(X) = -\sum_x p(x)\log_2 p(x),
  \label{eq:entropy}
\end{equation}
which is the Shannon entropy~\cite{CoverThomas}. It is the
minimum average bits per symbol achievable by any uniquely decodable
code. No compression algorithm can do better than $H(X)$ on average.

\subsection{Arithmetic Coding and the Model-Coder Separation}
\label{sec:arith}

Arithmetic coding solves the integer constraint of Huffman coding. The key property is that the coder is agnostic to the probability model.
Any sequential conditional estimator $\hat{p}(x_t \mid x_{1:t-1})$
plugs in directly. The total codelength is

\begin{equation}
  L(x^l) = \sum_{t=1}^{l} -\log_2 \hat{p}(x_t \mid x_{1:t-1}).
  \label{eq:codelength}
\end{equation}

With bits per symbol (BPS) calculated using

\begin{equation}
  \mathrm{BPS} = \frac{L(x^l)}{l}
  \label{eq:bps}
\end{equation}

When $\hat{p} = p_{\mathrm{true}}$, BPS equals the entropy rate. Any
deviation adds redundancy.

This separation is the mathematical engine of the paper. Compression
quality equals estimation quality. The transformer, Laplace, and CTW
are all probability estimators. Plugging any of them into~\eqref{eq:bps}
gives the codelength that arithmetic coding would achieve with that
model. BPS is computable directly from a code length calculator that uses this arithmetic coder, with no real compressor required. Every experiment in this paper reduces to: which estimator predicts test sequences better?

\subsection{Universal Source Coding}
\label{sec:universal}

A universal coder requires no knowledge of the source distribution. The redundancy of a code with length function $L$ for a source with
parameter $\boldsymbol{\theta}$ and sequences of length $l$ is
\begin{equation}
  R_l(L, \boldsymbol{\theta}) = \frac{1}{l}\mathbb{E}_{\boldsymbol{\theta}}[L(X^l)] - H_{\boldsymbol{\theta}}(X).
  \label{eq:redun}
\end{equation}
This is the per-symbol gap above entropy. 

This redundancy is minimized with a universal source coding algorithm that uses the \textit{Laplace estimator} to estimate bit probabilities sequentially. At each position $t$ it
estimates
\begin{equation}
  \hat{p}(x_{t+1}=1 \mid x_{1:t}) = \frac{k+1}{n+2},
  \label{eq:laplace}
\end{equation}
where $k$ is the count of ones seen and $n = t-1$ is the total symbols
seen. For bit prediction, this expression requires only running counts and no training.

CTW~\cite{gabeschamberg} blends Laplace estimates over all context lengths up to depth $d$, and is asymptotically more powerful than Laplace on sources with memory. Like Laplace, it requires no offline training, adapting instead to the test sequence during compression. To beat both CTW and Laplace estimation, a learned coder must exploit knowledge of the source distribution acquired through prior training.

%% ------------------------------------------------------------
\section{Experiment Design}
\label{sec:system}

\subsection{System Overview}

The full pipeline runs as follows. A source distribution generates
training sequences. The transformer is trained on those sequences, testing on sequences drawn from the same source. BPS is
computed by plugging transformer-generated bit predictions into
equation~\eqref{eq:bps}. Laplace and CTW are evaluated on the same
test sequences.

The paper reports six experiments. Experiment 1 trains a general
transformer on sequences drawn uniformly from $(0,1)$ and compares it
against Laplace. Experiment 2 trains specialized transformers per source
distribution and runs the same comparison. Experiment 2.5 trains
transformers under non-uniform Beta priors and compares against the
Bayes-optimal estimator for each prior. Experiment 3 sweeps the
training set size $m$ on a fixed IID source to measure redundancy
decay. Experiment 4 benchmarks a specialized transformer, the general IID
transformer, Laplace, and CTW on a binary symmetric Markov chain, showing
that the IID model's in-context gains are modest compared to specialized
training. Experiment 5 runs a Dirichlet sweep
over four-state Markov models, comparing a specialized transformer, the
general IID transformer, Laplace, and CTW. Experiment 6 benchmarks a
specialized transformer against Laplace and CTW on a four-state FSM source.

\subsection{The Transformer as Probability Estimator}
\label{sec:transformer}

The probability estimator is minGPT-nano~\cite{minGPT}, a scaled-down
decoder-only transformer with approximately $0.11$M parameters.
The architecture has three transformer blocks, each with three causal
self-attention heads and embedding dimension $d = 48$.
The vocabulary is binary ($V = 2$), with token embeddings and
learned positional embeddings both of dimension 48. The context length is
499 tokens, and therefore can be deployed for sequences up to 499 bits long.

The causal mask in each attention layer ensures that position $t$ attends
only to positions $1, \ldots, t$, enforcing that no future symbol is
visible. The input is the binary sequence $x_{1:t}$; the output at
position $t$ is a distribution over $\{0, 1\}$, from which we read off
$\hat{p}(x_{t+1} = 1 \mid x_{1:t})$. A single forward pass produces all
$l$ conditional probabilities simultaneously.

Training is done with learning rate $10^{-4}$ and batch size 64.
All models train for 1000 iterations on Apple MPS hardware.

\subsection{Arithmetic Source Coding and CTW}
\label{sec:fsm_def}

CTW and arithmetic coding form a complete, modular compression system. The
coder maintains an interval $[l, h) \subset [0, 1)$ and narrows it by
$\hat{p}(x_t \mid x_{1:t-1})$ at each step. CTW supplies those probabilities
by maintaining a binary tree of depth $d$, where each node holds a Laplace
estimate conditioned on the corresponding context suffix. The FSM source used
in Experiment~6 has deterministic transitions, so once CTW has seen enough
symbols its tree captures the state dynamics exactly. The transformer plugs
into the same arithmetic coding pipeline but gets its probabilities from a
neural network trained offline. The key difference is that the transformer has
already seen the FSM structure during training, while CTW must infer it from
scratch on every test sequence. 

%% ------------------------------------------------------------
\section{Experiments}
\label{sec:experiments}

The source parameter $p$ denotes the Bernoulli
bias. $l$ denotes test sequence length. BPS values are means over 100 trials;
error bars in all figures are one standard error.

%% -----------------------------------------------------------
\subsection{Experiment 1: The General Transformer Learns the Laplace Estimator}
\label{sec:exp1}

The first experiment trains a single ``general'' transformer on sequences
whose bias $p$ is drawn uniformly from $(0, 1)$ at
each training example. It sees the entire source distribution from
(0,1) simultaneously and does not accumulate a strong prior for any single
p. At test time it is evaluated
on sequences from (0, 1) and compared against Laplace.

\begin{figure}[t]
  \centering
  \includegraphics[width=\linewidth]{../notebooks/IID/figures/iid_bps_per_p.pdf}
  \caption{BPS vs.\ sequence length for the general transformer (blue) and
           Laplace (orange) across (0, 1). Entropy is shown as a
           dashed lower bound. Error bars are one standard error over 100 trials.}
  \label{fig:iid_per_p}
\end{figure}

Figure~\ref{fig:iid_per_p} shows the result. The general transformer tracks
Laplace across all five p values at every sequence length. At $p=0.3$, $l=10$,
Laplace achieves $0.980$ BPS and the general transformer achieves $0.990$ BPS.
By $l=499$ the gap narrows to $0.0005$ BPS. The pattern holds across the full
source distribution range.

The match is the finding. Laplace is Bayes-optimal under a uniform prior
on $p$, so a transformer that converges to the same predictions has internalized
an optimal Bayesian inference rule from data alone. The mixture training creates
a uniform effective prior, and the transformer's output distribution converges
to the corresponding optimal predictor.

\FloatBarrier
%% -----------------------------------------------------------
\subsection{Experiment 2: A Specialized Transformer Beats Laplace at Short Sequences}
\label{sec:exp2}

A specialized transformer is trained on sequences from a single fixed bias.
Two variants are evaluated: $p = 0.3$ and $p = 0.7$. Because training is
concentrated on one distribution, the model develops a strong implicit prior
for that source, similar in behavior to an optimal Bayesian estimator with a known prior.

\begin{figure}[t]
  \centering
  \includegraphics[width=0.49\linewidth]{../notebooks/IID/figures/iid_p03.png}
  \caption{Absolute BPS for
           specialized ($p=0.3$), general, Laplace, and entropy at $p=0.3$.}
  \label{fig:iid_gap}
\end{figure}

At $p=0.3$, $l=10$: the specialized transformer achieves $0.894$ BPS against
Laplace's $0.959$ BPS, a gap of $0.065$ bits per symbol. At $l=499$ the gap
collapses to $0.007$ BPS. The $p=0.7$ result is symmetric. This is because $H(0.3) = H(0.7)$ exactly: the binary entropy function is symmetric around $p=0.5$, so the two sources have identical entropy.

The specialized model's specialized
training acts as a useful prior that gives it an advantage when the
sequence is short. Laplace is
Bayes-optimal under a uniform prior, and as $l$ grows the empirical count
dominates any prior term, so Laplace catches up. The transformer's context
window is also capped at $499$ positions, so it cannot condition on arbitrarily
long histories the way a classical count-based estimator can.

\FloatBarrier
%% -----------------------------------------------------------
\subsection{Experiment 2.5: Prior Generalization to Beta}
\label{sec:exp_prior}

This experiment asks whether the specialized-beats-Laplace result is specific
to the uniform prior or generalizes. Separate transformers are trained under
five Beta priors: $\mathrm{Beta}(1,1)$, $\mathrm{Beta}(0.5,0.5)$,
$\mathrm{Beta}(5,5)$, $\mathrm{Beta}(5,2)$, and $\mathrm{Beta}(2,5)$.
Each model is evaluated against Laplace and a Bayes-optimal estimator
that uses the true prior analytically.

\begin{figure}[t]
  \centering
  \includegraphics[width=\linewidth]{../notebooks/ML/figures/fig1_bps_overview.pdf}
  \caption{BPS vs.\ sequence length for transformer, Laplace, CTW, and the
           Bayes-optimal estimator under each Beta prior. The entropy floor for
           each prior is shown as a dashed line.}
  \label{fig:prior_overview}
\end{figure}

Figure~\ref{fig:prior_overview} shows the result. Under $\mathrm{Beta}(1,1)$,
the transformer, Laplace, and estimator all align, consistent with Experiment~1.
CTW performs poorly at small $l$ across all non-uniform priors, never matching
the transformer at short sequences.

The transformer adapts to the prior shape and consistently sits close to the estimator for every prior tested. The gap to the estimator is small
but positive, which is expected since the transformer approximates the
Bayes-optimal rule rather than computing it exactly.

\FloatBarrier
%% -----------------------------------------------------------
\subsection{Experiment 3: Redundancy Scales as $1/m$}
\label{sec:scaling}

This experiment measures how transformer redundancy decays as training set
size $m$ grows. A fresh transformer is trained from scratch for each value
of $m \in \{10, 25, 50, 75, 100, 150, 200, 300, 500, 1000\}$ on IID sequences
with $p=0.3$. All models are tested on sequences of length $l=100$.
Redundancy is BPS minus $H(0.3) = 0.881$ bits/symbol.

\begin{figure}[t]
  \centering
  \includegraphics[width=0.8\linewidth]{../notebooks/train_bounds/figures/redundancy_vs_training.pdf}
  \caption{Transformer redundancy vs.\ training set size $m$ on a log-log scale.
           The Laplace redundancy (flat at $0.022$ bits/symbol) is shown as a
           dashed baseline. }
  \label{fig:scaling}
\end{figure}

Figure~\ref{fig:scaling} shows the result. At $m=10$: redundancy $0.373$
bits/symbol. At $m=50$: $0.040$. At $m=100$: $0.021$, already below the
Laplace baseline of $0.022$. By $m=150$: $0.012$, clearly dominant.
The crossover at $m \approx 100$ is where the pretrain-learned prior becomes
informative enough to outweigh Laplace's parameter-free advantage.

\FloatBarrier
%% -----------------------------------------------------------
\subsection{Experiment 4: Binary Symmetric Markov Chain}
\label{sec:markov}

This experiment runs four estimators on the binary symmetric Markov chain:
a specialized transformer trained per $p_{\mathrm{stay}}$, the general IID
transformer from Experiment~1, Laplace, and CTW (depth~8).

\begin{figure}[t]
  \centering
  \includegraphics[width=\linewidth]{../notebooks/Markov/figures/markov_bps_grid_comparison.png}
  \caption{BPS vs.\ sequence length on the binary symmetric Markov chain.
           Specialized transformer (dark blue), IID transformer (light blue),
           CTW depth-8 (brown), Laplace (green). At $p_{\mathrm{stay}}=0.5$
           the source is memoryless and all four estimators coincide.}
  \label{fig:markov_comparison}
\end{figure}

Figure~\ref{fig:markov_comparison} shows the result. The first observation
is that CTW and Laplace track closely across all $p_{\mathrm{stay}}$ values.
At $p_{\mathrm{stay}}=0.99$, $l=10$: Laplace $0.406$ BPS, IID transformer
$0.389$ BPS, CTW $1.000$ BPS. CTW has no full context window at $l=10$ with
depth~8, so it codes with flat $0.5$ priors throughout. By $l=50$ CTW
recovers ($0.346$ BPS) and pulls ahead of Laplace, and by $l=499$ CTW leads
while the IID transformer has fallen behind both universal coders.

The specialized transformer is a different story. Trained on sequences from
the same $p_{\mathrm{stay}}$, it learns the transition structure offline and
applies it immediately. At high $p_{\mathrm{stay}}$ it sits well below all
three other estimators at every sequence length, and the gap does not close
as $l$ grows.

The IID model's in-context gains over Laplace are real but small. CTW, a
universal coder with no offline training, matches or exceeds the IID
transformer at moderate to long sequence lengths. To consistently beat a universal coder on a structured source,
specialized training is required.

\FloatBarrier
%% -----------------------------------------------------------
\subsection{Experiment 5: Dirichlet Markov Model Sweep}
\label{sec:fsm_dirichlet}

This experiment tests a four-state Markov model with emission probabilities
$[0.1, 0.35, 0.65, 0.9]$ and a random transition matrix drawn per instance
from a $\mathrm{Dirichlet}(\alpha, \ldots, \alpha)$ distribution. Larger
$\alpha$ concentrates the transition matrix toward uniform (more randomness
per step); smaller $\alpha$ concentrates it toward sparser transitions.
Two sub-experiments are run for each
$\alpha \in \{0.1, 0.5, 1.0, 2.0, 5.0, 10.0\}$: one with a specialized
transformer trained per $\alpha$, and one applying the general IID transformer.

\begin{figure}[t]
  \centering
  \includegraphics[width=\linewidth]{../notebooks/Markov/figures/fsm_bps_grid_specialized.pdf}
  \caption{BPS vs.\ $l$ for the specialized transformer (blue), Laplace
           (orange), and CTW depth-8 (green) at six Dirichlet concentration
           parameters $\alpha$.}
  \label{fig:fsm_specialized}
\end{figure}

\begin{figure}[t]
  \centering
  \includegraphics[width=\linewidth]{../notebooks/Markov/figures/fsm_bps_grid_iid.pdf}
  \caption{Same comparison with the general IID transformer in place of the
           specialized model. The IID transformer does not beat Laplace or CTW.}
  \label{fig:fsm_iid}
\end{figure}

Figure~\ref{fig:fsm_specialized} shows the specialized transformer results.
At $\alpha=1.0$, $l=100$: transformer $0.994$ BPS, Laplace $1.009$ BPS, CTW
$1.022$ BPS. The specialized transformer wins at every $\alpha$ tested.

Figure~\ref{fig:fsm_iid} shows the IID transformer result. It does not beat
Laplace or CTW at any $\alpha$. This is consistent with Experiment~4: even
on the simpler binary symmetric Markov chain, the IID model's in-context
gains over universal coders were small and disappeared at longer sequences.
On a four-state Markov model with Dirichlet-random transitions, the IID
transformer offers nothing. Specialized training is required.

\FloatBarrier
%% -----------------------------------------------------------
\subsection{Experiment 6: FSM Source}
\label{sec:fsm_deterministic}

The FSM source provides the clearest
test of what offline training achieves against a proper universal coder. The three
estimators here are the specialized transformer, Laplace, and CTW at depth~8.
The FSM uses emission probabilities $[0.1, 0.35, 0.65, 0.9]$ and transition
matrices
\begin{equation}
T_0 = \begin{pmatrix}1 & 0 & 0 & 0 \\ 0 & 1 & 0 & 0 \\ 0 & 0 & 0 & 1 \\ 0 & 0 & 1 & 0\end{pmatrix},\quad
T_1 = \begin{pmatrix}0 & 1 & 0 & 0 \\ 0 & 0 & 1 & 0 \\ 0 & 0 & 0 & 1 \\ 1 & 0 & 0 & 0\end{pmatrix},
\label{eq:fsm_det}
\end{equation}
where the next state depends on both the current state and the emitted bit.

\begin{figure}[!t]
  \centering
  \includegraphics[width=0.85\linewidth]{../notebooks/Markov/figures/fsm_benchmark.pdf}
  \caption{BPS vs.\ $l$ for the specialized transformer, CTW depth-8, and
           Laplace on the four-state FSM source.}
  \label{fig:fsm_benchmark}
\end{figure}

Figure~\ref{fig:fsm_benchmark} shows the result. At $l=10$: transformer
$0.760$ BPS, Laplace $0.918$ BPS, CTW $1.000$ BPS. At $l=100$: transformer
$0.751$, CTW $0.840$, Laplace $0.827$. At $l=499$: transformer $0.747$,
CTW $0.798$, Laplace $0.810$.

The transformer leads at every sequence length. CTW closes the gap steadily
as $l$ grows due to online learning. At
short sequences the transformer's offline training is worth roughly $0.24$
BPS over CTW and $0.16$ BPS over Laplace. That gap is the operational value
of knowing the source distribution in advance.

%% -----------------------------------------------------------
\section{Discussion}
\label{sec:discussion}

The transformer is both a frozen coder and an in-context predictor at once. Its weights are fixed after offline training on $m$ sequences, and the distribution it applies to the test sequence is determined before any test bit arrives. On top of that, the transformer also reads the test sequence online, updating its prediction at each step. These two mechanisms are distinct, and the experiments show they contribute differently.

The pretrain-learned component is the prior encoded in the weights. When the transformer trains on sequences drawn from a specific source distribution, those weights encode a distribution over $p$ concentrated near the true distribution. This happens entirely offline before the test sequence exists. The in-context component is what the transformer does during coding: it conditions on the symbols it has already seen and refines its estimate of $p$ step by step. Both contribute to the final BPS, but they operate on different aspects and require different sets of training data.

The general transformer's convergence to Laplace demonstrates the in-context capability of transformers cleanly. The general transformer trains on a uniform distribution over $p$, which is exactly Laplace's assumed prior. Same prior, same posterior updates, same predictions. The two models agree because the pretrain-learned component contributes nothing beyond what Laplace already assumes. What remains is pure in-context learning, and on i.i.d.\ sequences that is enough to match Laplace but not quite beat it.

The specialized transformer adds a useful prior. Its weights arrive at the test sequence already concentrated near the right $p$. At short $l$, the in-context component has little data to work with, and this pretrain-learned head start dominates. As $l$ grows, empirical counts accumulate and the in-context component takes over for both the transformer and Laplace.

Experiment~2.5 shows the strongest version of this argument. The transformer does not just get a head start on Laplace. Trained under any Beta prior, it consistently tracks the Bayes-optimal estimator for that prior across all five families tested. This means the transformer is learning the Bayesian inference rule for the prior distribution. 

Both CTW and Laplace have no pretrain-learned component. They start from scratch on every test sequence and rely entirely on in-context adaptation. On structured sources at moderate to long sequences, CTW matches or beats the IID transformer. The IID transformer's pretrain-learned prior is not helpful on structured sources, and CTW's in-context reconstruction eventually catches up.

The specialized transformer beats CTW on structured sources, but the mechanism is the pretrain-learned prior, not superior in-context adaptation. CTW must reconstruct source structure from the test sequence from scratch, but the specialized transformer already knows it. On the FSM source that offline knowledge is worth roughly $0.24$ BPS over CTW at $l=10$.

%% -----------------------------------------------------------
\section{Conclusion}
\label{sec:conclusion}

A small transformer can learn to perform lossless source coding at near-optimal
rates. On IID sources, a general transformer recovers the Laplace estimator
from training data alone, and a specialized transformer beats Laplace at short
sequence lengths with the crossover around $m=100$ training sequences.
Redundancy decays at a rate consistent with the $1/m$ theoretical prediction.
On structured sources, the IID transformer shows small in-context gains over
Laplace on the binary symmetric Markov chain, but CTW achieves
better performance with no offline training. The gains from Laplace-only
comparisons do not hold against a proper universal coder baseline. Specialized
training is what actually beats CTW, and the advantage grows substantially on
four-state Markov and FSM sources.

The central takeaway is that learned compression and universal compression
are not competing approaches. They operate on the same mathematical engine of prediction.
The transformer earns its advantage by bringing offline knowledge of the
source distribution that classical universal coders cannot access.

%% -----------------------------------------------------------
\bibliographystyle{IEEEtran}
\bibliography{references}
\end{document}